\hypertarget{openbmc-code-update}{%
\section{OpenBMC Code Update}\label{openbmc-code-update}}

Two BMC Code layouts are available:

\begin{itemize}
\tightlist
\item
  Static, non-UBI layout
\item
  UBI layout - enabled via \texttt{obmc-ubi-fs} distro feature
\end{itemize}

This document describes the code update that supports both layouts.

\hypertarget{steps-to-update}{%
\subsection{Steps to Update}\label{steps-to-update}}

The following are the steps to update the BMC.

\begin{enumerate}
\def\labelenumi{\arabic{enumi}.}
\item
  Get a BMC image tar: After building OpenBMC, you will end up with a
  set of image files in
  \texttt{tmp/deploy/images/\textless{}platform\textgreater{}/}.

  \begin{itemize}
  \tightlist
  \item
    The UBI layout image is
    \texttt{obmc-phosphor-image-\textless{}platform\textgreater{}-\textless{}timestamp\textgreater{}.ubi.mtd.tar}
  \item
    The static layout image is
    \texttt{obmc-phosphor-image-\textless{}platform\textgreater{}-\textless{}timestamp\textgreater{}.static.mtd.tar}
  \end{itemize}

  The BMC tar image contains 5 files: u-boot, kernel, ro, and rw
  partitions and the MANIFEST file, which contains information about the
  image such as the image purpose, version, KeyType (Key type used for
  signature), HashType (SHA type used for key generation) and
  MachineName (name of machine used while building image, and this will
  be used for validation of image build). A MANIFEST file might look
  like

\begin{verbatim}
purpose=xyz.openbmc_project.Software.Version.VersionPurpose.BMC
version=2.7.0-dev
KeyType=OpenBMC
HashType=RSA-SHA256
MachineName=tiogapass
\end{verbatim}
\item
  Transfer the generated BMC image to the BMC via one of the following
  methods:

  \begin{itemize}
  \tightlist
  \item
    Method 1: Via Redfish Upload:
    \url{https://github.com/openbmc/docs/blob/master/REDFISH-cheatsheet.md\#firmware-update}.
    If using this method skip ahead to step 5!
  \item
    Method 2: Via scp: Copy the generated BMC image to the
    \texttt{/tmp/images/} directory on the BMC.
  \item
    Method 3: Via REST Upload:
    \url{https://github.com/openbmc/docs/blob/master/rest-api.md\#uploading-images}
  \item
    Method 4: Via TFTP: Perform a POST request to call the
    \texttt{DownloadViaTFTP} method of
    \texttt{/xyz/openbmc\_project/software}.
  \end{itemize}

  Methods 3 and 4 require additional options in bmcweb to be enabled.
\item
  Note the version id generated for that image file. The version id is a
  hash value of 8 hexadecimal numbers, generated by SHA-512 hashing the
  version string contained in the image and taking the first 8
  characters. Get the version id via one of the following methods:

  \begin{itemize}
  \tightlist
  \item
    Method 1: From the BMC command line, note the most recent directory
    name created under \texttt{/tmp/images/}, in this example it'd be
    \texttt{2a1022fe}:
  \end{itemize}

\begin{verbatim}
# ls -l /tmp/images/
  total 0
  drwx------    2 root     root            80 Aug 22 07:54 2a1022fe
  drwx------    2 root     root            80 Aug 22 07:53 488449a2
\end{verbatim}

  \begin{itemize}
  \tightlist
  \item
    Method 2: This method \emph{only} works if there are no
    \texttt{Ready} images at the start of transferring the image. Using
    the REST API, note the object that has its Activation property set
    to Ready, in this example it'd be \texttt{2a1022fe}:
  \end{itemize}

\begin{verbatim}
  $ curl -b cjar -k https://${bmc}/xyz/openbmc_project/software/enumerate
  {
    "data": {
      "/xyz/openbmc_project/software/2a1022fe": {
        "Activation": "xyz.openbmc_project.Software.Activation.Activations.Ready",
\end{verbatim}

  \begin{itemize}
  \tightlist
  \item
    Method 3: Calculate the version id beforehand from the image with:
  \end{itemize}

\begin{verbatim}
  tar xfO <BMC tar image> MANIFEST | sed -ne '/version=/ {s/version=//;p}' | head -n1 | tr -d '\n' | sha512sum | cut -b 1-8
\end{verbatim}
\item
  To initiate the update, set the \texttt{RequestedActivation} property
  of the desired image to \texttt{Active}, substitute
  \texttt{\textless{}id\textgreater{}} with the hash value noted on the
  previous step, this will write the contents of the image to the BMC
  chip via one of the following methods:

  \begin{itemize}
  \tightlist
  \item
    Method 1: From the BMC command line:
  \end{itemize}

\begin{verbatim}
  busctl set-property xyz.openbmc_project.Software.BMC.Updater \
    /xyz/openbmc_project/software/<id> \
    xyz.openbmc_project.Software.Activation RequestedActivation s \
    xyz.openbmc_project.Software.Activation.RequestedActivations.Active
\end{verbatim}

  \begin{itemize}
  \tightlist
  \item
    Method 2: Using the REST API:
  \end{itemize}

\begin{verbatim}
  curl -b cjar -k -H "Content-Type: application/json" -X PUT \
    -d '{"data":
    "xyz.openbmc_project.Software.Activation.RequestedActivations.Active"}' \
    https://${bmc}/xyz/openbmc_project/software/<id>/attr/RequestedActivation
\end{verbatim}
\item
  (Optional) Check the flash progress. This interface is only available
  during the activation progress and is not present once the activation
  is completed via one of the following:

  \begin{itemize}
  \tightlist
  \item
    Method 1: From Redfish: A task is returned from the Redfish upload.
    The task can be used to monitor the progress.
  \end{itemize}

\begin{verbatim}
 curl -k https://${bmc}/redfish/v1/TaskService/Tasks/0
\end{verbatim}

  \begin{itemize}
  \tightlist
  \item
    Method 2: From the BMC command line:
  \end{itemize}

\begin{verbatim}
  busctl get-property xyz.openbmc_project.Software.BMC.Updater  \
    /xyz/openbmc_project/software/<id> \
    xyz.openbmc_project.Software.ActivationProgress Progress
\end{verbatim}

  \begin{itemize}
  \tightlist
  \item
    Method 3: Using the REST API:
  \end{itemize}

\begin{verbatim}
  curl -b cjar -k https://${bmc}/xyz/openbmc_project/software/<id>/attr/Progress
\end{verbatim}
\item
  Check that the activation is complete by verifying the ``Activation''
  property is set to ``Active'' via one of the following methods:

  \begin{itemize}
  \tightlist
  \item
    Method 1: From Redfish: Check the task returned from the Redfish
    upload.
  \end{itemize}

\begin{verbatim}
  curl -k https://${bmc}/redfish/v1/TaskService/Tasks/0
\end{verbatim}

  \begin{itemize}
  \tightlist
  \item
    Method 2: From the BMC command line:
  \end{itemize}

\begin{verbatim}
  busctl get-property xyz.openbmc_project.Software.BMC.Updater \
    /xyz/openbmc_project/software/<id> \
    xyz.openbmc_project.Software.Activation Activation
\end{verbatim}

  \begin{itemize}
  \tightlist
  \item
    Method 3: Using the REST API:
  \end{itemize}

\begin{verbatim}
  curl -b cjar -k https://${bmc}/xyz/openbmc_project/software/<id>
\end{verbatim}
\item
  Reboot the BMC for the image to take effect.

  \begin{itemize}
  \item
    Method 1: From Redfish: If ApplyTime was set to ``Immediate'', the
    BMC will automatically reboot:
    \url{https://github.com/openbmc/docs/blob/master/REDFISH-cheatsheet.md\#firmware-applytime}.
    To reboot the BMC manually see:
    \url{https://github.com/openbmc/docs/blob/master/REDFISH-cheatsheet.md\#bmc-reboot}.
  \item
    Method 2: From the BMC command line:
  \end{itemize}

\begin{verbatim}
  reboot
\end{verbatim}

  \begin{itemize}
  \tightlist
  \item
    Method 3: Using the REST API:
  \end{itemize}

\begin{verbatim}
  curl -c cjar -b cjar -k -H "Content-Type: application/json" -X PUT \
      -d '{"data": "xyz.openbmc_project.State.BMC.Transition.Reboot"}' \
      https://${bmc}/xyz/openbmc_project/state/bmc0/attr/RequestedBMCTransition
\end{verbatim}
\end{enumerate}

\hypertarget{associations}{%
\subsection{Associations}\label{associations}}

In addition to all software images, several associations are listed at
\texttt{/xyz/openbmc\_project/software/}:

\begin{verbatim}
curl -c cjar -b cjar -k -H "Content-Type: application/json" -X GET  \
    https://${bmc}/xyz/openbmc_project/software/
{
  "data": [
    "/xyz/openbmc_project/software/46e65782",
    "/xyz/openbmc_project/software/493a00ad",
    "/xyz/openbmc_project/software/88c153b1",
    "/xyz/openbmc_project/software/active",
    "/xyz/openbmc_project/software/functional"
  ],
  "message": "200 OK",
  "status": "ok"
}
\end{verbatim}

\begin{enumerate}
\def\labelenumi{\arabic{enumi}.}
\item
  A ``functional'' association to the ``running'' BMC and host images

  There is only one functional association per BMC and one functional
  association per host. The functional/running BMC image is the BMC
  image with the lowest priority when rebooting the BMC. The functional
  image does not update until the BMC is rebooted. The functional host
  image behaves the same way except that it updates on a power on or
  reboot of the host.

\begin{verbatim}
curl -c cjar -b cjar -k -H "Content-Type: application/json" -X GET \
    https://${bmc}/xyz/openbmc_project/software/functional
{
  "data": {
    "endpoints": [
      "/xyz/openbmc_project/software/46e65782",
      "/xyz/openbmc_project/software/493a00ad"
    ]
  },
  "message": "200 OK",
  "status": "ok"
}
\end{verbatim}
\item
  An ``active'' association to the active BMC and host images

  Note: Several BMC images might be active, this is true for the host
  images as well.

\begin{verbatim}
curl -c cjar -b cjar -k -H "Content-Type: application/json" -X GET \
    https://${bmc}/xyz/openbmc_project/software/active
{
  "data": {
    "endpoints": [
      "/xyz/openbmc_project/software/46e65782",
      "/xyz/openbmc_project/software/493a00ad",
      "/xyz/openbmc_project/software/88c153b1"
    ]
  },
  "message": "200 OK",
  "status": "ok"
}
\end{verbatim}
\item
  An ``updateable'' association to the programmable components

  This is used for identifying firmware components which are
  programmable via BMC OOB interfaces like Redfish/IPMI. All updateable
  firmware components must expose the updateable association so that
  upper applications like Redfish/IPMI will know about updateable
  firmwares.

  To know the updateable software components:

\begin{verbatim}
 busctl call xyz.openbmc_project.ObjectMapper \
  /xyz/openbmc_project/software/updatable org.freedesktop.DBus.Properties \
  Get ss xyz.openbmc_project.Association endpoints
v as 1 "/xyz/openbmc_project/software/1201fc36"
\end{verbatim}

  Redfish interface uses `updateable' association in SoftwareInventory
  schema.
\item
  An additional association is located at
  \texttt{/xyz/openbmc\_project/software/\textless{}id\textgreater{}/inventory}
  for ``associating'' a software image with an inventory item.

\begin{verbatim}
curl -c cjar -b cjar -k -H "Content-Type: application/json" -X GET \
   https://${bmc}/xyz/openbmc_project/software/493a00ad/inventory
{
  "data": {
    "endpoints": [
      "/xyz/openbmc_project/inventory/system/chassis/motherboard/boxelder/bmc"
    ]
  },
  "message": "200 OK",
  "status": "ok"
}
\end{verbatim}

  To get all software images associated with an inventory item:

\begin{verbatim}
curl -c cjar -b cjar -k -H "Content-Type: application/json" -X GET  \
    https://${bmc}/xyz/openbmc_project/inventory/system/chassis/activation
{
  "data": {
    "endpoints": [
      "/xyz/openbmc_project/software/46e65782"
    ]
  },
  "message": "200 OK",
  "status": "ok"
}
\end{verbatim}
\end{enumerate}

\hypertarget{manifest-file}{%
\subsection{MANIFEST File}\label{manifest-file}}

A file named ``MANIFEST'' must be included in any image tar uploaded,
downloaded via TFTP, or copied to the BMC.

The MANIFEST file format must be key=value (e.g.~version=v1.99.10). It
should include the following fields:

\begin{itemize}
\tightlist
\item
  version - The version of the image
\item
  purpose - The image's purpose (e.g.
  xyz.openbmc\_project.Software.Version.VersionPurpose.BMC or
  xyz.openbmc\_project.Software.Version.VersionPurpose.Host). Accepted
  purpose values can be found at
  \href{https://github.com/openbmc/phosphor-dbus-interfaces/blob/6f69ae5b33ee224358cb4c2061f4ad44c6b36d70/xyz/openbmc_project/Software/Version.interface.yaml}{Version
  interface} under ``VersionPurpose'' values.
\item
  MachineName - The name of machine (platform) for which this image is
  built for. This value will be compared against
  OPENBMC\_TARGET\_MACHINE value defined in os-release file of running
  image. Image will not be upgraded if this check fails. For backward
  compatibility this check skips failure if MachineName is not defined
  for current released images but it will be made mandatory field from
  2.9 onward releases.
\end{itemize}

Other optional fields are:

\begin{itemize}
\tightlist
\item
  extended\_version - A more detailed version, which could include
  versions of different components in the image.
\end{itemize}

\hypertarget{deleting-an-image}{%
\subsection{Deleting an Image}\label{deleting-an-image}}

To delete an image:

\begin{verbatim}
curl -c cjar -b cjar -k -H "Content-Type: application/json" \
    -X POST https://${bmc}/xyz/openbmc_project/software/<$id>/action/delete \
    -d "{\"data\": [] }"
\end{verbatim}

Note: The image must be non-functional (``non-running'').

To delete all non-functional images, whether BMC or host images:

\begin{verbatim}
curl -c cjar -b cjar -k -H "Content-Type: application/json" \
    -X POST https://${bmc}/xyz/openbmc_project/software/action/deleteAll \
    -d "{\"data\": [] }"
\end{verbatim}

\hypertarget{software-field-mode}{%
\subsection{Software Field Mode}\label{software-field-mode}}

Field mode is meant for systems shipped from manufacturing to a
customer. Field mode offers a way to provide security and ensure
incorrect patches don't get loaded on the system by accident. The
software implementation of the field mode interface disables patching of
the BMC by not mounting \texttt{/usr/local}, which in turn disables host
patching at \texttt{/usr/local/share/pnor/}. Enabling field mode is
intended to be a one-way operation which means that once enabled, there
is no REST API provided to disable it.

Field mode can be enabled by running the following command:

\begin{verbatim}
curl -b cjar -k -H 'Content-Type: application/json' -X PUT -d '{"data":1}'  \
    https://${bmc}/xyz/openbmc_project/software/attr/FieldModeEnabled
\end{verbatim}

Although field mode is meant to be a one-way operation, it can be
disabled by a user with admin privileges by running the following
commands on the BMC:

\begin{verbatim}
fw_setenv fieldmode

systemctl unmask usr-local.mount

reboot
\end{verbatim}

More information on field mode can be found here:
\url{https://github.com/openbmc/phosphor-dbus-interfaces/blob/master/yaml/xyz/openbmc_project/Control/FieldMode.interface.yaml}

\hypertarget{software-factory-reset}{%
\subsection{Software Factory Reset}\label{software-factory-reset}}

Software factory reset resets the BMC and host firmware to its factory
state by clearing out any read/write data. To software factory reset run
the following command and then reboot the BMC:

\begin{verbatim}
curl -b cjar -k -H 'Content-Type: application/json' -X POST -d '{"data":[]}' \
    https://${bmc}/xyz/openbmc_project/software/action/Reset
\end{verbatim}

The factory reset on the BMC side will clear \texttt{/var},
\texttt{/home}, and \texttt{/etc}. On the host side, the factory reset
will clear the read/write volume for each host image on the system,
clear the shared preserve host volume, pnor-prsv, and clear any host
patches located in \texttt{/usr/local/share/pnor/}.

The factory reset interface can be found here:
\url{https://github.com/openbmc/phosphor-dbus-interfaces/blob/02b39246d45ea029a1652a49cc20eab7723dd63b/xyz/openbmc_project/Common/FactoryReset.interface.yaml}

\hypertarget{image-storage-location}{%
\subsection{Image Storage Location}\label{image-storage-location}}

\hypertarget{static-layout}{%
\subsubsection{Static layout}\label{static-layout}}

When a BMC image is activated, each
\texttt{image-\textless{}name\textgreater{}} is written to the BMC
chip's partitions indicated by the
\texttt{\textless{}name\textgreater{}}:

\begin{itemize}
\tightlist
\item
  image-u-boot
\item
  image-kernel
\item
  image-rofs
\item
  image-rwfs
\end{itemize}

\hypertarget{ubi-layout}{%
\subsubsection{UBI layout}\label{ubi-layout}}

When a BMC image is activated (i.e.~when ``RequestedActivation'' is set
to ``Active''), UBI volumes are created on the BMC chip for the image.
The alternate BMC chip can also be used to store images. This is
determined by ``BMC\_RO\_MTD''. Using both the alternate BMC chip and
the BMC chip allows for multiple BMC images to be stored. By default,
only the BMC chip is used. To use both, set ``BMC\_RO\_MTD'' to
``alt-bmc+bmc''.

\hypertarget{implementation}{%
\subsection{Implementation}\label{implementation}}

More information about the implementation of the code update can be
found at
\url{https://github.com/openbmc/phosphor-dbus-interfaces/blob/master/yaml/xyz/openbmc_project/Software}
and \url{https://github.com/openbmc/phosphor-bmc-code-mgmt}
